\documentclass[12pt]{article}
\usepackage[utf8]{inputenc}
\usepackage[margin=1in]{geometry}
\usepackage{amsmath}
\usepackage{graphicx}
\usepackage{booktabs}
\usepackage{hyperref}
\usepackage[round]{natbib}

\title{Nicotine-Related Beliefs Induce Dose-Dependent Responses in a Thalamic Circuit in Human Smokers}
\author{Anonymous Authors}
\date{}

\begin{document}
\maketitle

\begin{abstract}
Beliefs about drug potency can powerfully modulate subjective experience and neural function, yet whether such modulation follows a graded, dose-dependent relationship---analogous to pharmacological dose-response curves---remains unknown. Here we tested whether instructed beliefs about e-cigarette nicotine strength produce parametric, dose-dependent changes in brain activity during value-based decision-making. Twenty nicotine-dependent adults completed three fMRI sessions, each paired with a different belief condition (``low,'' ``medium,'' or ``high'' nicotine), while actual nicotine content was held constant at 1.2\%. Beliefs successfully modulated perceived nicotine strength (rated 3.52, 4.52, 5.82 on a 10-point scale; $F(2,38) = 9.71$, $P = 0.0004$, partial $\eta^2 = 0.34$), while actual nicotine intake, metabolism, and baseline saturation did not differ across conditions. Whole-brain analysis revealed that thalamic activity during value tracking increased parametrically with belief dosage (cluster-level $P_{\text{FWE}} = 0.006$), consistent with the thalamus's high density of nicotinic acetylcholine receptors. Functional coupling between thalamus and ventromedial prefrontal cortex also scaled parametrically to belief dosage. In contrast, striatal reward-related responses (nucleus accumbens) did not differ across belief conditions ($P = 0.945$), demonstrating anatomical specificity. These findings reveal that subjective beliefs about nicotine strength can produce graded, dose-dependent neural responses in circuits known to be pharmacologically sensitive to nicotine, establishing a bridge between cognitive constructs and neurochemical dose-response relationships.
\end{abstract}

\section{Introduction}

The influence of beliefs and expectations on neural function has been a central focus of placebo research \citep{wager2004placebo, benedetti2005neurobiological, colloca2013placebo}. Placebo effects have been demonstrated across diverse domains including pain, motor function, and reward processing, typically as binary comparisons between active and placebo conditions. However, pharmacological agents produce dose-dependent effects: neural responses to nicotine, alcohol, and other substances scale with the administered dose. Whether beliefs alone can produce similarly graded, dose-dependent neural modulation has never been tested.

Nicotine is a particularly suitable system for investigating belief-driven neural dose-response relationships. The thalamus contains one of the highest densities of nicotinic acetylcholine receptors (nAChRs) in the human brain, and PET imaging studies have demonstrated that thalamic nAChR occupancy responds to nicotine in a dose-dependent manner \citep{cosgrove2009imaging, brody2004thalamic}. Importantly, habitual smoking behavior can produce substantial nAChR occupancy even with minimal nicotine, suggesting that expectation and habit may contribute to receptor-level responses independently of the pharmacological agent \citep{picciotto2012acetylcholine}.

Prior work from our group demonstrated that binary belief manipulations (believing nicotine is present vs.\ absent) modulate reward-related neural signals in smokers \citep{gu2015belief}. However, this binary design cannot distinguish between an all-or-none belief effect and a graded dose-response relationship. Establishing graded belief-driven neural responses would fundamentally advance our understanding of how cognitive constructs interface with neurochemical systems, with implications for addiction, placebo medicine, and computational psychiatry \citep{friston2014computational}.

Here we employed a within-subject design in which nicotine-dependent smokers completed three fMRI sessions, each paired with a different instructed belief about nicotine strength (low, medium, or high), while actual nicotine content was held constant. We focused on the thalamus and striatum as regions of interest based on their known pharmacological sensitivity to nicotine and their roles in reward processing and value-based decision-making \citep{kable2009subjective, brody2006nicotine}. A healthy non-smoking control cohort provided an exploratory baseline for thalamic and striatal reward-tracking signals in the absence of nicotine exposure.

The study design incorporates multiple safeguards to ensure that observed neural effects can be attributed to beliefs rather than pharmacology. First, the within-subject design controls for all stable individual differences in brain structure, receptor density, and baseline function. Second, multiple biological measures---nicotine intake via cartridge weight, cotinine via mass spectrometry, and exhaled carbon monoxide---independently confirm that pharmacological exposure does not differ across belief conditions. Third, the three-level belief manipulation (low, medium, high) enables testing for parametric dose-response relationships that would be impossible with a simple binary placebo/active design.

\section{Methods}

\subsection{Participants}

Twenty-three nicotine-dependent adults were enrolled; three were excluded (1 software malfunction, 1 fell asleep during scanning, 1 lost behavioral data), yielding a final sample of 20 participants (6 female, 14 male; mean age $41.1 \pm 11.97$ years, range 24--61; all right-handed). Inclusion criteria required smoking a minimum of 10 cigarettes per day for over 1 year, with no prior e-cigarette use, no current quit attempts, and no illicit drug use in the past 2 months (Table~\ref{tab:participants}).

A separate healthy control cohort of 31 non-smokers (15 female, 16 male; mean age $28 \pm 9$ years; 2 excluded for excessive head movement $> 3$~mm from initial 33) completed the same decision-making task without nicotine manipulation for exploratory comparison.

Power analysis based on a prior nicotine-belief fMRI study (Cohen's $d = 0.69$) indicated a minimum sample of 18--20 participants for 90\% power at $\alpha = 0.05$.

\begin{table}[ht]
\centering
\caption{Participant characteristics.}
\label{tab:participants}
\begin{tabular}{lcc}
\toprule
Parameter & Smoker Cohort & Control Cohort \\
\midrule
Enrolled & 23 & 33 \\
Final N & 20 & 31 \\
Sex (F/M) & 6/14 & 15/16 \\
Age (mean $\pm$ SD) & $41.1 \pm 11.97$ & $28 \pm 9$ \\
Age range & 24--61 & -- \\
Handedness & All right & -- \\
Sessions & 3 & 1 \\
\bottomrule
\end{tabular}
\end{table}

\subsection{Experimental Design}

Each participant completed three fMRI sessions spaced approximately one week apart. In each session, participants were instructed that an e-cigarette contained either ``low,'' ``medium,'' or ``high'' nicotine (order counterbalanced). Actual nicotine content was always 1.2\%. The session proceeded as follows: (1) pre-vaping saliva collection for cotinine measurement, (2) exhaled CO measurement as baseline saturation index, (3) belief instruction, (4) 20 minutes of controlled vaping, (5) post-vaping saliva collection, and (6) fMRI scanning during a value-based decision-making task.

\subsection{Manipulation Checks}

Perceived nicotine strength was rated on a 10-point scale after each session. Nicotine intake was computed as the change in e-cigarette cartridge weight multiplied by 1.2\%. Cotinine concentration was quantified in saliva samples via high-performance liquid chromatography tandem mass spectrometry (LC-MS/MS). Exhaled CO was measured before vaping.

\subsection{fMRI Analysis}

The primary analysis examined reward-tracking signals (parametric modulation by market return) across the three belief conditions using a repeated-measures ANOVA design. A priori regions of interest included the thalamus (independent anatomical mask) and nucleus accumbens. Psychophysiological interaction (PPI) analysis \citep{friston1997psychophysiological} assessed belief-dependent changes in functional connectivity between thalamus and ventromedial prefrontal cortex (vmPFC). Family-wise error correction was applied at the cluster level ($k = 50$).

\section{Results}

\subsection{Manipulation Checks}

The belief manipulation was successful: perceived nicotine strength increased parametrically across conditions (Table~\ref{tab:checks}). Critically, actual nicotine intake, cotinine changes, and exhaled CO levels did not differ across conditions, confirming that any neural differences between belief conditions could not be attributed to differences in pharmacological exposure.

\begin{table}[ht]
\centering
\caption{Summary of manipulation checks across belief conditions.}
\label{tab:checks}
\begin{tabular}{lccccc}
\toprule
Measure & Low & Medium & High & $F$ statistic & $P$ \\
\midrule
Perceived strength (AU) & 3.52 & 4.52 & 5.82 & $F(2,38) = 9.71$ & 0.0004 \\
Nicotine intake (mg) & 0.928 & 0.719 & 0.783 & $F(2,38) = 1.806$ & 0.178 \\
Cotinine change (ng/mL) & -- & -- & -- & $F(2,32) = 0.959$ & 0.393 \\
Exhaled CO (ppm) & -- & -- & -- & $F(2,32) = 0.364$ & 0.698 \\
\bottomrule
\end{tabular}
\end{table}

\subsection{Thalamic Dose-Dependent Response}

Whole-brain analysis revealed a significant cluster in the thalamus where reward-tracking activity increased parametrically with belief dosage (cluster-level $P_{\text{FWE}} = 0.006$). ROI analysis using an independent thalamic mask confirmed this parametric increase ($P = 0.036$), with the lowest activity under the low-belief condition, intermediate under medium belief, and highest under high belief. Across individuals, subjective perception of nicotine strength correlated parametrically with thalamic neural activity, bridging subjective experience and neural representation.

\subsection{Striatal Specificity}

In stark contrast to the thalamus, the nucleus accumbens showed no parametric modulation by belief condition. While striatal reward-tracking signals were present across all conditions when collapsed (consistent with engagement by the decision-making task), the rmANOVA comparing belief conditions was far from significant ($P = 0.945$; permutation test $P = 0.94$). This double dissociation---present in thalamus, absent in striatum---demonstrates the anatomical specificity of the belief-driven dose-response effect (Table~\ref{tab:regions}).

\begin{table}[ht]
\centering
\caption{Comparison of belief-driven modulation in thalamus and striatum.}
\label{tab:regions}
\begin{tabular}{llll}
\toprule
Region & Dose-Response & $P$ (rmANOVA) & Interpretation \\
\midrule
Thalamus & Present (parametric) & 0.036 (ROI) & Belief-sensitive \\
Nucleus accumbens & Absent & 0.945 & Belief-insensitive \\
Thalamus--vmPFC coupling & Present (parametric) & Significant & Circuit-level effect \\
\bottomrule
\end{tabular}
\end{table}

\subsection{Thalamus--vmPFC Functional Connectivity}

PPI analysis revealed that functional coupling between the thalamus and vmPFC also scaled parametrically with belief dosage. This circuit-level dose-response relationship indicates that beliefs about nicotine strength modulate not only regional thalamic activity but also the communication between the thalamus and a cortical region implicated in value representation and state inference.

\section{Discussion}

We demonstrate that instructed beliefs about nicotine strength produce graded, dose-dependent neural responses in the thalamus of nicotine-dependent smokers, with no corresponding modulation in the striatum. This finding extends prior binary placebo-nicotine research \citep{gu2015belief} by showing that the relationship between beliefs and neural activity follows a parametric dose-response function analogous to pharmacological dose-response curves.

The anatomical specificity of the effect---present in thalamus but absent in nucleus accumbens---is theoretically meaningful. The thalamus contains one of the highest densities of nAChRs in the human brain \citep{cosgrove2009imaging, picciotto2012acetylcholine}, making it a natural candidate for belief-driven modulation that mimics pharmacological effects. The posterior thalamus is particularly enriched in nAChRs and plays a central role in gating incoming sensory information and enhancing attention, functions that are directly affected by nicotine. That beliefs can parametrically modulate activity in this region, without corresponding modulation in the striatum, suggests that belief-driven neural effects may preferentially target circuits with high receptor density for the believed substance.

The parametric increase in thalamus--vmPFC coupling with belief dosage reveals a circuit-level mechanism through which beliefs may exert their neural effects. The vmPFC is broadly implicated in value computation, state representation, and model-based inference \citep{kable2009subjective, friston2014computational}. The strengthened coupling between thalamus and vmPFC under higher belief conditions may reflect a top-down mechanism through which belief-derived state inferences modulate thalamic processing.

The robust dose-dependent subjective ratings (partial $\eta^2 = 0.34$) confirm that participants genuinely experienced the three belief conditions as qualitatively different levels of nicotine exposure, despite identical pharmacological content. This subjective grading, combined with the neural grading, establishes a chain linking instructed belief to subjective experience to thalamic activity to thalamocortical connectivity.

Several limitations warrant consideration. The sample size of 20 smokers, while adequately powered for the primary analysis, limits the ability to examine individual differences in belief susceptibility. The e-cigarette delivery method may introduce variability in nicotine absorption that is not fully captured by cartridge weight changes. The healthy control comparison is exploratory and confounded by age differences between cohorts. Finally, the within-subject design, while powerful for controlling individual differences, may be susceptible to carry-over effects across the approximately one-week inter-session intervals.

\section{Conclusion}

Beliefs about nicotine strength produce graded, dose-dependent neural responses in thalamic circuits that are known to be pharmacologically sensitive to nicotine, with anatomical specificity that spares the striatum. These findings establish that cognitive constructs can produce parametric neural effects that parallel pharmacological dose-response relationships, providing a neurobiological foundation for understanding how beliefs shape drug experience and suggesting new targets for belief-based interventions in addiction treatment.

\bibliographystyle{plainnat}
\bibliography{references}

\end{document}
